\input{template.tex}

\ifPDFTeX
  \renewcommand{\rmdefault}{cmr}
\fi

\def\SheetNo{6}
\setlength{\emergencystretch}{2em}

\begin{document}
\sheettitle{Разбиения + комбинаторика.}

\noindent В задачах о разбиениях числа на натуральные слагаемые порядок слагаемых не учитывается. Разбиение нуля - пустая сумма; разбиений отрицательного числа нет.

\begin{problems}

\item На доске написано несколько натуральных чисел $a_0,a_1,\ldots,a_n$. На второй доске записывают числа: $b_0$ - сколько всего чисел на первой доске, $b_1$ - сколько там чисел, больших единицы, $b_2$ - сколько чисел, больших двойки, и так далее, пока получаются положительные числа. Нули не записывают. На третьей доске записывают числа $c_0,c_1,\ldots$, построенные по числам второй доски по тому же правилу. Докажите, что наборы чисел на первой и третьей досках совпадают с учётом кратностей; порядок чисел не важен.

\item В классе учатся $32$ человека. Для участия в стритболе они организовали $33$ различные команды по три человека в каждой. Докажите, что найдутся две команды, у которых ровно один общий участник.

\item Вершины правильного стоугольника покрашены в $10$ цветов. Докажите, что можно выбрать четыре его вершины, являющиеся вершинами прямоугольника и покрашенные не более чем в два различных цвета.

\item Тест по английскому языку сдавали $10$ школьников. Известно, что любые пять школьников вместе ответили на все вопросы, а любые четыре вместе ответили не на все вопросы. При каком наименьшем количестве вопросов теста такое могло случиться? Группа ответила на вопрос, если на него ответил хотя бы один участник группы.

\item Пусть $0\le n\le k\ell$. Докажите, что количество разбиений числа $n$ в сумму не более чем $k$ слагаемых, каждое из которых не превосходит $\ell$, равно количеству разбиений числа $k\ell-n$ в сумму не более чем $k$ слагаемых, каждое из которых не превосходит $\ell$.

\item Петя написал количество разбиений числа $n$ в сумму не более чем $k$ слагаемых, каждое из которых не превосходит $\ell$. Вася написал количество разбиений числа $n$ в сумму не более чем $k$ слагаемых, каждое из которых не превосходит $\ell-1$. Маша написала количество разбиений числа $n-\ell$ в сумму не более чем $k-1$ слагаемых, каждое из которых не превосходит $\ell$. Докажите, что сумма Машиного и Васиного чисел равна числу Пети.

\item Три пирата Джо, Билл и Том нашли клад, содержащий $80$ одинаковых золотых монет, и хотят разделить их так, чтобы каждому досталось не менее $15$ монет. Сколько существует способов это сделать?

\item Петя подсчитал количество всех возможных $m$-буквенных слов из букв $\mathrm{T}$, $\mathrm{O}$, $\mathrm{W}$ и $\mathrm{N}$, в каждом из которых букв $\mathrm{T}$ и $\mathrm{O}$ поровну. Вася подсчитал количество всех возможных $2m$-буквенных слов из букв $\mathrm{T}$ и $\mathrm{O}$, в каждом из которых этих букв поровну. У кого слов получилось больше? Слово - любая последовательность букв.

\item Найдите количество восьмизначных чисел, произведение цифр которых равно $1400$. Ответ представьте в виде целого числа.

\item Паша, Оля и Саша играли в тетрис. Победитель очередного раунда получал одно очко, остальные - по нулю. Игра заканчивается, как только кто-то наберёт $101$ очко. У Оли всё ещё было $0$ очков, когда один из двух остальных участников впервые набрал $100$. Но затем Оля собралась с силами и победила, причём в конце игры у Паши и Саши было по $100$ очков. Сколько существует последовательностей побед в раундах, при которых такое могло произойти?

\item Дано натуральное $n$. Докажите, что числа от $1$ до $2^n$ можно покрасить в два цвета так, чтобы не было одноцветной арифметической прогрессии длины $2n$. Разность прогрессии должна быть положительной.

\item Даны $6000$ карточек, на которых написаны натуральные числа от $1$ до $6000$, каждое ровно один раз. Сколькими способами можно выбрать две карточки, чтобы сумма написанных на них чисел делилась на $100$? Порядок выбранных карточек не учитывается.

\item Сколько семибуквенных слов можно составить из букв $a$ и $b$ так, чтобы непосредственно после каждой буквы $a$ стояла хотя бы одна буква $b$? Букве $a$ разрешается быть последней.

\item Петя хочет выписать все последовательности из $100$ натуральных чисел, в каждой из которых хотя бы раз встречается число $4$ или $5$, а любые два соседних члена отличаются не больше чем на $2$. Сколько последовательностей ему придётся выписать?

\item Сколькими способами можно представить $10^{10}$ в виде произведения трёх натуральных множителей, если произведения, отличающиеся порядком множителей, считаются различными?

\item Сколько существует $23$-значных чисел, сумма цифр которых равна восьми?

\item Изначально все клетки доски $8\times8$ белые. За один ход можно поменять на противоположный цвет все клетки любого квадрата $3\times3$, все клетки любого квадрата $4\times4$ или одну из клеток $a1$ и $b2$. Любую ли чёрно-белую раскраску доски можно получить такими ходами? Обозначения клеток - обычные шахматные.

\item Школьники участвовали в олимпиаде, проходившей в два тура. В каждый из двух дней они рассаживались по нескольким кабинетам, причём никто не сидел в кабинете в одиночестве. Количество кабинетов в разные дни могло отличаться. Докажите, что найдутся два школьника $A$ и $B$ такие, что в первый день они писали олимпиаду в кабинетах с одинаковым количеством участников и во второй день также писали в кабинетах с одинаковым количеством участников. Сами кабинеты у этих школьников могли быть разными.

\item Докажите, что количество разбиений числа $n$ на натуральные слагаемые равно количеству разбиений числа $2n$ на ровно $n$ натуральных слагаемых.

\item Число называется \emph{шикарным}, если у него есть два делителя, отличающихся на $2$, и \emph{фееричным}, если у него есть два делителя, отличающихся на $3$. Каких чисел больше от $1$ до $6000$ включительно - шикарных или фееричных?

\item Докажите, что при любом натуральном $n$ уравнения
\[
x^2+y^2=n
\qquad\text{и}\qquad
x^2+y^2=2n
\]
имеют одинаковое количество решений в целых числах.

\item Рассмотрим всевозможные наборы из $100$ неотрицательных целых чисел, расположенных в неубывающем порядке и не превосходящих $100$, в которых сумма всех чисел кратна $10$. Докажите, что ровно половина этих наборов заканчивается на $100$.

\item Докажите, что количество разбиений числа $n$ в сумму не более чем $k$ слагаемых, каждое из которых не превосходит $\ell$, равно количеству разбиений числа $n$ в сумму не более чем $\ell$ слагаемых, каждое из которых не превосходит $k$.

\item Сколько существует способов разбить число $2015$ на приблизительно равные натуральные слагаемые? Слагаемых может быть одно или несколько. Числа называются приблизительно равными, если их разность не больше $1$. Например:
\[
\underbrace{1+1+\ldots+1}_{2015\text{ слагаемых}},
\qquad 671+672+672,
\qquad 1007+1008.
\]

\item Назовём билет с номером от $000000$ до $999999$ отличным, если некоторые две соседние цифры его номера отличаются на $5$. Найдите число отличных билетов. Нули в начале номера допускаются.

\end{problems}
\end{document}
